\chapter{Resultados e Discussão}
\label{capitulo:resultados_discussao}

Esta seção é dinâmica e dependerá muito do seu trabalho.
Mas nela podem ser incluídas as principais telas desenvolvidas no sistema.
Discussões sobre diferenças entre o protótipo e a tela.

Preferencialmente peça para seu usuário julgar as telas do seu sistema.
Para isto elabore uma metodologia, que deve ser incluída no capítulo \ref{capitulo:materiais_metodos}, de avaliação do sistema desenvolvido.

Nesta metodologia não precisa ser complexa, basta colocar uma ficha contendo notas para aparência, facilidade de uso, tempo de resposta, utilidade, etc. Facilite sua análise pedindo para o usuário dar notas de 0 a 10 ou em outra escala, o importante é que seja numérica.

Debata sobre estes resultados nesta seção. Observe que se seu trabalho é desenvolver um novo algoritmo e os critérios de avaliação são claros, não é necessário que um usuário o julgue.
Veja, você está desenvolvendo um novo algoritmo de ordenação, a complexidade e medições de do tempo gasto na ordenação já bastarão\footnote{Note que isto não evita que a metodologia de análise seja definida no capítulo \ref{capitulo:materiais_metodos}.}.
